\section{System Description}
\label{sec:system}

This section describes the 11-stage pipeline as a designed system: its stage structure,
artifact dependencies, gate conditions, AI autonomy model, and implementation substrate.
The system contribution is the pipeline specification itself --- a reusable process
architecture that can be instantiated across hunt domains without modification.

% ----------------------------------------------------------------
\subsection{Pipeline Overview}
\label{sec:pipeline-overview}

The pipeline consists of eleven sequential stages, each producing a structured artifact
that serves as the primary input to the next stage. Figure~\ref{fig:pipeline} provides a
schematic of the full stage sequence; the two review gates (Stages~3 and~7) are marked as
decision points. Table~\ref{tab:stages} provides the full stage reference.

\begin{table}[h]
\centering
\caption{Pipeline stage reference. AI autonomy levels: \textbf{FA} = fully AI-autonomous;
\textbf{AP} = AI-primary, human-confirms; \textbf{PD} = AI-panel-driven (review gate);
\textbf{HP} = human-primary. Gate condition applies only to stages with a formal pass/fail
threshold.}
\label{tab:stages}
\small
\begin{tabular}{clllll}
\toprule
\textbf{\#} & \textbf{Stage} & \textbf{Primary Input} & \textbf{Primary Output} &
\textbf{Gate Condition} & \textbf{AI Role} \\
\midrule
1 & Scope         & Hunt brief (prose)           & \texttt{SCOPE.md}, \texttt{CLAUDE.md}       & Human approval      & AP \\
2 & World Build   & \texttt{SCOPE.md}            & \texttt{world/}, \texttt{LOCKED.md}         & ---                 & AP \\
3 & Puzzle Pool   & World data, \texttt{ROUNDS.md} & \texttt{PUZZLE-POOL.md} (scored)           & Score $\geq$ 33/45  & PD \\
4 & Assignment    & \texttt{PUZZLE-POOL.md}      & Per-puzzle \texttt{ASSIGNMENT.md}           & ---                 & FA \\
5 & Meta Design   & Feeder answers               & \texttt{META-DESIGN.md}                     & Human verification  & AP \\
6 & Authoring     & Per-puzzle \texttt{ASSIGNMENT.md} & Puzzle files (one per puzzle)          & Score $\geq$ 22/30  & FA \\
7 & Editorial     & Puzzle files                 & Per-puzzle \texttt{REVIEW.md}               & Pass/revise/cut     & PD \\
8 & Integration   & Editorial-passed puzzles     & \texttt{PUZZLES.md} integration checklist   & ---                 & FA \\
9 & Delivery Build & Integrated puzzle set        & Site, print, or props package               & ---                 & AP \\
10 & Platform Test & Delivery package             & Solver simulation report                    & ---                 & FA \\
11 & Polish        & Simulation report            & Final artifacts, hint system                & ---                 & HP \\
\bottomrule
\end{tabular}
\end{table}

The pipeline is \emph{artifact-driven}: each stage's output is a version-controlled
markdown document, and subsequent stages treat those documents as read-only inputs. This
constraint enforces unidirectional information flow and prevents coherence drift, which
would otherwise occur when late-stage authoring decisions propagate backwards to modify
scope or world documents. A dedicated \emph{world lock} enforces this at the content level:
Stage~2 produces a \texttt{LOCKED.md} canon declaration, and all subsequent stages are
prohibited from introducing facts not present in the locked world files. Every factual
claim in every puzzle must be traceable to a specific entry in the \texttt{world/}
directory. This mechanism was validated empirically: across nine hunt scenarios, no
post-lock canon violations were detected in Stage~7 editorial review, compared to a
baseline of multiple coherence errors per scenario in early pilot runs conducted without
the lock mechanism.

Stage~1 (Scope) establishes the hunt's identity: domain, audience, scale, tone, delivery
format, and thematic frame. The primary output is a \texttt{SCOPE.md} document and a
per-scenario \texttt{CLAUDE.md} status tracker. The status tracker is the pipeline's
runtime state: each stage updates it on completion, and the \texttt{/hunt resume} skill
reads it to recover from session interruptions without duplicating completed work.

Stage~2 (World Build) populates the content library from which puzzles draw their
facts. For fictional domains (Ironfall's 1993 JRPG universe, Dead Reckoning's
starship systems), the world is synthesized and populated into structured data tables.
For real-world domains (Age of Empires, Wavelength's pop music corpus), the world
stage collects, verifies, and canonicalizes factual data. The \texttt{/hunt world}
skill organizes this data into individual system files (bestiary, items, locations,
characters, and so on) and then issues the canon lock.

Stage~3 (Puzzle Pool) generates a pool of candidates --- typically three times the
number of puzzles needed for the final hunt --- and subjects them to AI panel ranking
(see Section~\ref{sec:gates}). The scored pool is the primary input to Stage~4
(Assignment), which allocates the highest-scoring candidates to puzzle slots and
coordinates all feeder answer words to ensure meta compatibility.

Stage~5 (Meta Design) is intentionally ordered before Stage~6 (Authoring): the meta
puzzle's extraction mechanism determines which feeder answer words are required, and
all feeder puzzles must be briefed with their target answer words before authoring
begins. This ordering prevents the most common failure mode in AI-generated hunts ---
meta answers that do not match the feeder set --- by establishing answer contracts
before any puzzle content is written.

Stage~6 (Authoring) writes puzzle files from the coordinated briefs. In the
multi-scenario runs described in Section~\ref{sec:empirical}, Stage~6 was executed as
parallel AI agent launches, with one or more independent agents per puzzle, running
simultaneously without shared state. Each agent reads its assignment brief, the relevant
world data files, and the design principles document, then produces a complete puzzle
page. No human editing occurs between brief and authored output.

Stages~7 through~11 constitute the review-to-ship pipeline: editorial review (with gate),
cross-puzzle integration verification, delivery format construction, platform simulation,
and final polish. Stage~11 is the only stage coded as human-primary (see
Section~\ref{sec:ai-roles}): it encompasses hint writing, edge case detection, and final
answer encoding --- tasks that require human judgment about solver experience that AI
simulation does not reliably approximate.

% ----------------------------------------------------------------
\subsection{Gate Conditions}
\label{sec:gates}

The pipeline contains two primary review gates: Stage~3 (pool ranking) and Stage~7
(editorial review). Both are implemented using the AI expert panel system described
in \cite{dellaLibera2025panel}. A secondary gate operates within Stage~6 (authoring):
each puzzle must pass an automated quality threshold before proceeding to editorial.

\paragraph{Stage~3: Pool Ranking Gate.}

The puzzle pool gate filters candidates before the investment of authoring work. The panel
evaluates each candidate against the seven-dimension weighted rubric introduced in
\cite{dellaLibera2025panel}: Clarity ($\times 1$), Solvability ($\times 1$), Elegance
($\times 2$), Reading Reward ($\times 2$), Fun ($\times 1$), Confirmation ($\times 1$),
and Riven Standard ($\times 1$), yielding a maximum score of 45. At Stage~3, candidates
are evaluated on their \emph{brief} --- the one-to-two paragraph specification of the
puzzle concept, mechanism, and extraction --- rather than on authored content. The pass
threshold is 33/45 (73\%), matching the threshold calibrated in the rubric validation
experiment of \cite{dellaLibera2025panel} against three quality tiers of published puzzle
hunt puzzles.

Candidates scoring below 33 are not assigned. Human override is permitted: a hunt
designer may promote a candidate below threshold if the brief's weakness is judged
recoverable in authoring, or suppress a high-scoring candidate for editorial reasons
(e.g., domain overlap with another selected puzzle). In practice, across six scenarios,
human override occurred in fewer than 5\% of pool decisions; the panel's ranking was
treated as authoritative in the large majority of cases.

The Stage~3 gate is designed to fail cheaply. At the point of pool evaluation, the only
work invested is the candidate brief (typically 200--400 words). Failing a candidate at
Stage~3 therefore wastes at most the brief --- not the full authoring effort, which for a
complete puzzle is an order of magnitude larger. The gate is explicitly calibrated to
have high recall rather than high precision: it is better to promote a borderline candidate
that fails at Stage~7 than to suppress a strong candidate that a brief undersells.

\paragraph{Stage~7: Editorial Review Gate.}

The editorial gate operates on fully authored puzzles and produces a verdict in one of
three categories: \textsc{pass} (puzzle proceeds to integration), \textsc{revise} (puzzle
returns to Stage~6 with a specific revision brief), or \textsc{cut} (puzzle is removed
from the hunt). The panel evaluates authored content against the same rubric used in
Stage~3, augmented by three additional editorial checks: extraction chain verification
(character-by-character), voice consistency audit (does the puzzle match the hunt's
established narrative voice), and world-lock compliance (does every factual claim
trace to a locked world file).

A \textsc{revise} verdict triggers a targeted revision pass: the authoring agent receives
the review feedback and rewrites the specific extraction mechanism, clueing logic, or
presentation element flagged by the panel. In the scenarios reported here, revisions were
executed in a single pass in all cases where revision was required; no puzzle required
more than one revision cycle. A \textsc{cut} verdict was reserved for puzzles with
fundamental mechanism failures that would require redesign from the brief, which was judged
more efficient than revision.

The Stage~7 gate is the pipeline's primary quality control mechanism. Unlike the Stage~3
gate, which screens design concepts, Stage~7 screens executed construction. The failure
modes it catches are different: not ``is this concept worth pursuing'' but ``does this
puzzle actually work as built.'' The empirical data (Section~\ref{sec:empirical}) shows
that Stage~7 catches failures that Stage~3 does not predict: concept-level scores are
not reliable predictors of construction-level failures, which are typically mechanism
errors (extraction chains that do not resolve) or answer coordination errors (feeder
answers that do not match meta requirements).

\paragraph{Why These Two Gates.}

The two-gate design reflects a deliberate cost-staging argument. Creative pipeline failures
are not equally expensive at all stages; the cost of failure is proportional to the
downstream investment that must be discarded. Stage~3 failures are cheap (brief only);
Stage~7 failures are moderate (authoring effort only, before integration); failures that
reach Stage~8 or beyond are expensive (integration, delivery build, and test work is lost
or must be redone). Placing review gates at Stages~3 and~7 maximizes the expected quality
of work that reaches integration while minimizing total evaluation overhead. No gate is
placed between Stage~7 and Stage~11 because the editorial and integration stages are
jointly sufficient to catch remaining failures --- which, empirically, are concentrated
in delivery format issues rather than puzzle construction failures.

% ----------------------------------------------------------------
\subsection{AI Roles by Stage}
\label{sec:ai-roles}

Each stage in the pipeline is coded on a four-level autonomy gradient: \textbf{FA} (fully
AI-autonomous), \textbf{AP} (AI-primary, human confirms), \textbf{PD}
(AI-panel-driven, with human deciding on revisions), and \textbf{HP} (human-primary).
Figure~\ref{fig:autonomy} illustrates the gradient across the full pipeline. The coding
is based on the empirical record of how each stage was executed across all scenarios
in this study.

\emph{Stages~2, 4, 6, 8, and~10} are fully AI-autonomous in practice. Stage~2 (World
Build) generates and populates content tables from domain specifications with no human
editing of individual entries. Stage~4 (Assignment) translates pool rankings to puzzle
slot assignments and generates all briefs without human intervention. Stage~8 (Integration)
runs the integration checklist --- answer verification, cross-reference checking, round
structure validation --- autonomously. Stage~10 (Platform Test) executes a solver
simulation against the delivery package, with each simulated solver personality working
through the hunt independently and logging failures.

Stage~6 (Authoring) warrants specific discussion as the most autonomous stage in the
pipeline. Puzzles are authored by independent AI agents launched in parallel, one per
puzzle (or one per module in large hunts), with no human review occurring between brief
and authored output. The agents share no state; they communicate only through the
coordinated answer words established in Stage~5. In the largest scenario in this
study (Cosmere, 36 feeder puzzles across four modules), 36 puzzles were authored in
parallel across a single session without human editorial intervention until Stage~7. This
represents a qualitative shift from human-assisted authoring: the human's role in Stage~6
is to launch agents and wait for completion, not to review or approve individual puzzle
content. Stage~6 autonomy is feasible precisely because Stage~5 has pre-coordinated all
answer words and Stage~7 will catch authoring failures.

\emph{Stage~1 (Scope) and Stage~5 (Meta Design)} are coded AP: AI generates the scope
document or meta mechanism, and the human confirms before proceeding. These confirmations
are lightweight --- typically a review of the SCOPE.md document or a meta extraction trace
--- but they are genuine decision points at which the human can redirect the hunt's
identity before downstream work commits to a particular direction.

\emph{Stages~3 and~7} are panel-driven. The AI panel produces scored rankings (Stage~3)
or pass/revise/cut verdicts (Stage~7), but the human decides on override and revision
scope. In practice, these stages are AI-primary in execution and human-primary in
decision authority; the panel output is treated as a strong recommendation, not an
automatic action.

\emph{Stage~11 (Polish)} is the pipeline's most human-intensive stage. Hint writing
requires calibrating the difficulty of each hint tier to the expected solver population
--- a judgment about cognitive load that AI simulation does not reliably approximate. Edge
case detection requires the human to reason about degenerate solver paths that the platform
test simulation may not have exercised. Final answer encoding and security checks (ensuring
no plaintext answers appear in any distributed file) are procedurally defined but require
human sign-off before the hunt ships. Stage~11 is also where the hunt designer's creative
voice is most directly expressed: the final pass shapes the solver's experience in ways
that are not captured by the quality rubric.

% ----------------------------------------------------------------
\subsection{Implementation}
\label{sec:implementation}

The pipeline is implemented as a \emph{CLAUDE.md skill library}: a collection of
persistent markdown documents, each encoding a named skill as a structured prompt with
defined inputs, defined outputs, and explicit instructions for interacting with the
scenario's artifact directory. Skills are installed into Claude Code's persistent skill
store (the \texttt{\textasciitilde{}/.claude/skills/} directory) and invoked by slash
command (e.g., \texttt{/hunt world}, \texttt{/puzzle P01 author}). No bespoke code is
written for any scenario; the entire pipeline executes through natural language skill
invocations against version-controlled artifact files.

Skills are organized into two namespaces. The \texttt{/hunt} namespace contains
hunt-level administrative skills run by the hunt director: \texttt{plan} (the full
pipeline walkthrough), \texttt{world} (world build and lock), \texttt{review} (panel
review invocation), \texttt{meta} (meta design and verification), \texttt{edit}
(editorial review pass), \texttt{resume} (crash recovery), \texttt{print} (print package
generation), \texttt{site} (website scaffold and puzzle pages), \texttt{props} (physical
prop logistics), \texttt{publish} (solver-facing distribution package), \texttt{script}
(game architecture script for interactive delivery), \texttt{status} (pipeline
dashboard), and \texttt{profile} (reviewer profile management). The \texttt{/puzzle}
namespace contains puzzle-level skills used by authoring agents: \texttt{author}
(write puzzle from brief), \texttt{test} (blind test against quality threshold), and
\texttt{check} (design principles audit).

The skill library grew from 0 to 13 skills over the two-day study period (February~27--28,
2026), with skills added reactively in response to pipeline needs discovered during
scenario execution. Table~\ref{tab:skills} in Section~\ref{sec:empirical} provides the
full growth timeline. This organic development trajectory is itself an empirical finding:
the 13 skills that emerged are those that were actually needed to run the pipeline end-to-end,
not a pre-specified design. Skills added late in the study (notably \texttt{/hunt script}
and \texttt{/hunt publish}) address delivery and distribution tasks that only become
visible as a bottleneck when multiple scenarios reach Stage~9 simultaneously.

Each scenario's per-scenario \texttt{CLAUDE.md} references the toolkit skill library
and specifies only three per-scenario parameters: domain (the content universe or real-world
topic), audience (target solver population and experience level), and scale (puzzle count,
round structure, and delivery format). No per-scenario modifications to skill prompts were
required across any of the nine scenarios in this study. The same \texttt{/hunt world lock},
\texttt{/puzzle P01 author}, and \texttt{/hunt edit} commands that ran the Age of Empires
scenario in February also ran the Cosmere scenario with 36 puzzles and four meta layers,
without any prompt engineering between scenarios. This cross-theme invariance is the
primary evidence for the pipeline's transferability claim.

The \texttt{/hunt resume} skill deserves specific mention as an infrastructure component.
Claude Code sessions are stateful but not persistent: a session crash loses the in-context
working memory of all stages executed in that session. The resume skill recovers from
crashes by reading the per-scenario \texttt{CLAUDE.md} status tracker (updated at every
stage boundary), identifying the first incomplete stage, and regenerating the necessary
context from the artifact files on disk before continuing. In Stage~6, resume behavior is
puzzle-aware: it checks each puzzle individually and re-authors only puzzles that do not
yet have a completed file, skipping puzzles that were successfully authored before the
crash. Across the multi-session, multi-scenario runs of this study, the resume skill was
invoked at least once in the majority of scenarios, confirming that crash recovery is a
practical operational requirement rather than an edge case.

\begin{figure}[h]
\centering
% Placeholder for TikZ pipeline diagram
% Figure 1: 11-stage pipeline with gate markers at Stages 3 and 7,
% input/output artifact labels, and autonomy coding per stage.
\fbox{\parbox{0.9\textwidth}{\centering [Figure~1: Pipeline diagram --- 11 stages in
sequence, with gate condition markers at Stages~3 and~7, input/output artifact labels
per stage, and autonomy gradient coding. TikZ source in \texttt{figures/fig1-pipeline.tex}.]}}
\caption{The 11-stage AI-assisted puzzle hunt pipeline. Review gates (Stages~3 and~7)
are the primary quality control interventions. The autonomy gradient shifts from
AI-primary in the middle stages (authoring, integration, testing) to human-primary at
the boundaries (scope and polish).}
\label{fig:pipeline}
\end{figure}
